
\include{./def/yf-formatting}
\include{./def/yf-def}

\usepackage{wasysym}
%\usepackage{times}

\title{}
\author{}
\date{}


%\def\A{\mathcal{A}}
%\def\B{\mathcal{B}}
\def\C{\mathcal{C}}
%\def\D{\mathcal{D}}
\def\E{\mathcal{E}}
\def\FF{\EuScript{F}}
\def\G{\mathcal{G}}
\def\GG{\EuScript{G}}
\def\H{\mathcal{H}}
\def\I{\mathcal{I}}
\def\J{\mathcal{J}}
\def\L{\mathcal{L}}
\def\Q{\mathcal{Q}}
\def\S{\mathcal{S}}
\def\T{\mathcal{T}}
\def\V{\mathcal{V}}
%\def\W{\EuScript{W}}
\def\W{\mathscr{W}}
\def\X{\mathcal{X}}
\def\Y{\mathcal{Y}}
\def\Z{\mathcal{Z}}
\def\att{\mit{diffatt}}
\def\agm{\mathit{AGM}}
\def\boldatt{\mathbf{att}}
\def\ext{\mathit{ext}}
\def\In{\mathrm{IN}}
\def\join{\mathit{join}}
\def\Out{\mathrm{OUT}}
\def\schema{\mathit{schema}}

\def\extraspacing{\vspace{4mm} \noindent}
\def\vgap{\vspace{4mm}}
\def\vslit{\vspace{2mm}}

\begin{document}

\boxminipg{\linewidth}{
    \begin{center}
        \vspace{3mm}
    {\Large
    (CE7456 Lecture Notes 14) Massively Parallel Computation} \\[2mm]
    {\large Yufei Tao}
    \vspace{3mm}
    \end{center}
}

We have seen two computational models so far: (i) the RAM model, which is suitable when the data fit in main memory, and (ii) the EM model, which is suitable when the data fit on the disk (of one machine). In practice, many applications need to process datasets so large that they must be stored on multiple machines. The performance bottleneck in processing such data --- e.g., in systems like Hadoop and Spark --- is neither CPU computation nor I/O access; rather, it is typically the communication across machines. The RAM and EM models are no longer appropriate for studying algorithms in such scenarios.

\vgap

This lecture introduces the {\em massively parallel computation} (MPC) model \cite{bks17,ksv10} for studying algorithms designed to minimize parallel communication over a large number of machines. We will then illustrate the characteristics of the model through several representative algorithms.

\extraspacing {\bf Math Conventions.} Given an integer $x \ge 1$, we use $[x]$ to represent the set $\set{1, 2, ..., x}$.

\section{The MPC Model} \label{sec:model}

In this model, we have $p$ share-nothing machines that are interconnected in a network. Denote by $S$ the input set of elements, and let $n = |S|$. In the beginning, each machine stores $O(n/p)$ elements of $S$; an algorithm has no control over how elements are distributed across the machines. As a standard assumption, the values of $p$ and $n$ satisfy $p < n^{1-\eps}$, where $\eps$ is a positive constant less than 1.


\vgap

An algorithm starts by having each machine perform some initial computation on its local data. Then, it executes in {\em rounds}, each having two phases:
\myitems{
  \item in the first phase, the machines exchange messages (every message should have been prepared either in the initial computation or the second phase of the previous round);
  \item in the second phase, each machine performs local computation.
}
An algorithm can perform any number of rounds as needed. Upon termination, the final output must be available on the machines in the form specified by the problem definition.

\vgap

The {\em load of a round} is the largest number of words communicated (sent and received combined) by a machine in that round; specifically, if machine $i \in [p]$ communicates $\ell_i$ words in this round, then the load of the round is $\max_{i=1}^p \ell_i$.
The (total) {\em load} of the algorithm is the sum of the loads of all the rounds.

\vgap

With load $n$, any problem can be solved trivially in one round by simply sending all data to one machine. In practice, a good algorithm ought to perform a small number of rounds and incur a small load. Accordingly, two interesting questions arise in theory:
\myitems{
    \item ({\bf Load complexity}) What is the smallest load for solving a problem within a given number of rounds?

    \item ({\bf Round complexity}) What is the smallest number of rounds for solving a problem when each round can incur no more than a given load?
}

\section{Sorting} \label{sec:sort}

The sorting problem is defined as follows in the MPC model. Let $S$ be a set of $n$ integers, initially distributed across $p$ machines in an arbitrary manner, subject to the constraint that each machine has $O(n/p)$ elements. The goal is to re-distribute the elements across the machines such that the elements on machine $i$ are less than those on machine $j$, for any $1 \le i < j \le p$.

\vgap

\todo{State the guarantees}

\extraspacing {\bf Algorithm.} For each $i \in [p]$, denote by $S_i$ the set of elements stored on machine $i$ initially.

\vgap

In the first phase of round 1, each machine generates a sample set $\Sigma_i \subseteq S_i$ as follows: for each element of $S_i$, independently toss a coin with probability $\rho$ of heads ---  where $\rho$ is a value to be determined later --- and add the element to $\Sigma_i$ if the coin comes up heads.

\vgap

In the second phase of round 1, machine $i \in [p]$ sends $\Sigma_i$ to every other machine.

\vgap

In the first phase of round 2, each machine collects
\myeqn{
    S_\mit{sam} &=& \bigcup_{j=1}^p \Sigma_j. \nn
}
Set
\myeqn{
    s = |S_\mit{sam}|. \nn
}
The machine identifies $p - 1$ elements $b_1, b_2, ..., b_{p-1}$ where $b_j$ ($1 \le j \le p-1$) is the $j \cdot \lc s/p \rc$-th smallest elements in $S_\mit{sam}$. In addition, define $b_0 = -\infty$ and $b_p = \infty$. We refer to $b_0, b_1, ..., b_p$ as {\em boundary elements}.

\vgap

In the second phase of round 2, machine $i \in [p]$ sends to machine $j \in [p]$ the elements in $S_i \cap (b_{i-1}, b_i]$.

\extraspacing {\bf Analysis.} Define
\myeqn{
    m = n/p. \nn
}
The algorithm incurs a load of $O(m + p \ln n)$ when three conditions are met (think: why?):
\myitems{
    \item {{\bf C1:}} $|\Sigma_i| = O(\ln n)$ for all $i \in [p]$.
    %\item {{\bf C2:}} $s = O(m)$.
	\item {{\bf C2:}} $S \cap (b_{i-1}, b_i]$ has $O(m)$ elements for all $i \in [p]$.
}

\vgap


\begin{theorem} \label{thm:sort-tera-samplerate}
	When $m \ge t \ln (nt)$, {\bf C1} and {\bf C2} hold simultaneously with probability at least $1-1/n^2$ by setting $\rho = \fr{1}{m} \ln(nt)$.
\end{theorem}

\begin{proof}
	We will consider $p \ge 9$ because otherwise $m = \Omega(n)$, in which case {\bf C1} and {\bf C2} hold trivially. Our proof is based on the Chernoff bound and a bucketing argument.

	\vgap

	Consider an arbitrary $i \in [p]$. Let random variable $X$ be the size of $\Sigma_i$. Hence, $\expt[X] = m \rho = \ln(nt)$. An application of Chernoff bound \eqref{eqn:chernoff2} gives:
	\begin{eqnarray}
		\Pr[X \ge 18 \ln(nt)] \le \exp(- 3 \ln(nt)) \le 1/n^3. \nonumber
	\end{eqnarray}
	By union bound, the probability that this is true for at least one $i \in [p]$ is at most $p/n^3$. It thus follows that {\bf C1} holds with probability at least $1 - p/n^3$.

	\vgap

	Next, we consider {\bf C2}.
	Imagine that $S$ has been sorted in ascending order. We divide the sorted list into $\lf p/8 \rf$ sub-lists as evenly as possible; call each sub-list a {\em bucket}. Each bucket has between $8n/t=8m$ and $16m$ objects. We observe that {\em {\bf C2} holds if every bucket covers at least one boundary object}. To understand why, notice that under this condition, no bucket can fall between two consecutive boundary objects (counting also the dummy ones)\footnote{If there was one, the bucket would not be able to cover any boundary object.}. Hence, every $S_i$, $1 \le i \le t$, can contain objects in at most 2 buckets, i.e., $|S_i| \le 32m = O(m)$.

	\vgap

	A bucket $\beta$ definitely includes a boundary object if $\beta$ covers more than $1.6 \ln(nt) > s/t$ samples (i.e., objects from $S_\mit{sam}$), as a boundary object is taken every $\lc s/t \rc$ consecutive samples. Let $|\beta| \ge 8m$ be the number of objects in $\beta$. Define random variable $x_j$, $1 \le j \le |\beta|$, to be 1 if the $j$-th object in $\beta$ is sampled, and 0 otherwise. Define:
	\begin{eqnarray}
		X = \sum_{j=1}^{|\beta|} x_j = |\beta \cap S_\mit{sam}|. \nn
	\end{eqnarray}
	Clearly, $\expt[X] \ge 8m\rho = 8 \ln(nt)$. We have:
	\begin{eqnarray}
		\Pr[X \le 1.6 \ln(nt)] &=& \Pr[X \le (1 - 4/5)8 \ln(nt)] \nn \\
		&\le& \Pr[X \le (1 - 4/5)\expt[X]] \nn \\
		\textrm{(by Chernoff)} &\le& \exp\left(- \fr{16}{25} \fr{\expt[X]}{3}\right) \nn \\
		&\le& \exp\left(- \fr{16}{25} \cdot \fr{8 \ln(nt)}{3}\right) \nn \\
		&\le& \exp(-\ln(nt)) \nn \\
		&\le& 1/(nt). \nn
	\end{eqnarray}
	We say that $\beta$ {\em fails} if it covers no boundary object. The above derivation shows that $\beta$ fails with probability at most $1/(nt)$. As there are at most $t/8$ buckets, the probability that at least one bucket fails is at most $1/(8n)$. Hence, {\bf C2} can be violated with probability at most $1/(8n)$ under the event $s < 1.6 t \ln(nt)$, i.e., at most $9/8n$ overall.

	\vgap

	Therefore, {\bf C1} and {\bf C2} hold at the same time with probability at least $1 - 17/(8n)$.
\end{proof}


%\section{Remark} ksv10

\bibliographystyle{plainurl}% the mandatory bibstyle
\bibliography{ref}

\appendix

\section*{Appendix: Chernoff Bounds}

%\begin{lemma} \label{lmm::chernoff}
    Let $X_1, X_2, ..., X_f$ be independent Bernoulli random variables that are identically distributed. Define $X = \sum_{i \in [f]} X_i$. For any $\gamma \geq 2$, it holds that
    \myeqn{
        \Pr[X \geq \gamma \cdot \expt[X]] &\leq&  \exp(-\gamma \cdot \expt[X] / 6). \label{eqn:chernoff1}
    }
    For any $0 < \gamma < 1$, it holds that
    \myeqn{
        \Pr[X \geq (1 + \gamma) \expt[X]] &\leq&  \exp(- \gamma^2 \cdot \expt[X] / 3). \label{eqn:chernoff2}
    }
%\end{lemma}
Proofs for \eqref{eqn:chernoff1} and \eqref{eqn:chernoff2} can be found in \cite{stl12}.

\end{document}
